\subsection{半范数与凸集}

设 \(\mathbf{E}\) 是 \(\mathbf{R}\) 上的向量空间；映射 \(p\)（\(\mathbf{E}\) 到 \(\mathbf{R}\) 中的）称为正齐次的，如果对每个 \(\lambda \geqslant 0\) 及所有 \(x \in \mathbf{E}\)，我们有

\[
p (\lambda x) = \lambda p (x).\tag{5}
\]

E 上的正齐次函数 p 是凸的，当且仅当它对 E 的所有 x、y 满足 II, 第 1 页 的公理 \((\mathrm{SN}_{\mathrm{II}})\)：

\[
p (x + y) \leqslant p (x) + p (y).\tag{6}
\]

事实上，若 \(p\) 是凸的，则对 \(x, y\)（\(E\) 中的），

\[
p \left(\frac {1}{2} (x + y)\right) \leqslant \frac {1}{2} p (x) + \frac {1}{2} p (y)
\]

且由 (5)，这个关系式等价于 (6)。反之，若 (6) 成立，则对所有 \(\lambda\)（满足 \(0 < \lambda < 1\)），我们也有

\[
p (\lambda x + (1 - \lambda) y) \leqslant p (\lambda x) + p ((1 - \lambda) y) = \lambda p (x) + (1 - \lambda) p (y)
\]

这里利用了 (5)。

E 上的凸正齐次函数称为次线性函数。

若 \(p\) 是定义在 \(\mathbf{E}\) 上的次线性函数，则由 II, § 2.8, corollary，对所有 \(a > 0\)，集 \(\mathrm{V}(p, a)\)（相应地 \(\mathrm{W}(p, a)\)）——由点 \(x \in \mathrm{E}\)（满足 \(p(x) \leqslant a\)（相应地 \(p(x) < a\)））组成——是凸的；进而这个集是吸收的，因为对所有 \(x \in \mathrm{E}\)，存在 \(\lambda > 0\) 使得 \(p(\lambda x) = \lambda p(x) < a\)。

这个结果有一个部分逆命题：

命题 22. --- 设 A 是向量空间 E 中包含 0 的一个凸集。对所有 \(x \in \mathbf{E}\)，置

\[
p _ {\mathrm{A}} (x) = \inf _ {\rho > 0, x \in \rho \mathrm{A}} \rho\tag{7}
\]

\((0 \leqslant p_{\mathbf{A}}(x) \leqslant \infty)\)。函数 \(p_{\mathbf{A}}\) 满足

\[
p _ {\mathrm{A}} (x + y) \leqslant p _ {\mathrm{A}} (x) + p _ {\mathrm{A}} (y), \quad p _ {\mathrm{A}} (\lambda x) = \lambda p _ {\mathrm{A}} (x)\tag{8}
\]

对所有 \(x, y\)（\(\mathbf{E}\) 中的）及 \(\lambda > 0\) 成立。若用 \(\mathrm{V}(p_{\mathbf{A}}, \alpha)\)（相应地 \(\mathrm{W}(p_{\mathbf{A}}, \alpha)\)）表示诸 \(x \in \mathbf{E}\)（满足 \(p_{\mathbf{A}}(x) \leqslant \alpha\)（相应地 \(p_{\mathbf{A}}(x) < \alpha\)））组成的集，则

\[
\mathrm{W} (p _ {\mathrm{A}}, 1) \subset \mathrm{A} \subset \mathrm{V} (p _ {\mathrm{A}}, 1).\tag{9}
\]

若 \(\mathbf{A}\) 是吸收的，则 \(p_{\mathbf{A}}\) 是有限的（从而是次线性的）。

由于关系式 \(x \in \rho A\) 与 \(\lambda x \in \lambda \rho A\) 当 \(\lambda > 0\) 时等价，我们有 \(p_A(\lambda x) = \lambda p_A(x)\)（对 \(\lambda > 0\)）。设 \(x, y\) 是 \(E\) 的两点。若 \(x\)（相应地 \(y\)）不被 \(A\) 吸收，则 \(p_A(x) = +\infty\)（相应地 \(p_A(y) = +\infty\)），而不等式 \(p_A(x + y) \leqslant p_A(x) + p_A(y)\) 显然成立。设存在 \(\alpha > 0\)、\(\beta > 0\) 使得 \(x \in \alpha A\) 且 \(y \in \beta A\)；则 \(x + y \in \alpha A + \beta A = (\alpha + \beta)A\)（II, 第 8 页, 注记）；从而 \(p_A(x + y) \leqslant p_A(x) + p_A(y)\)。包含关系 \(A \subset V(p_A, 1)\) 显然成立。包含关系 \(W(p_A, 1) \subset A\) 则由 \(A\) 是凸的且包含 0 得出。最后，若 \(A\) 是吸收的，则 \(p_A\) 显然是有限的。

由 (7) 定义的函数 \(p_{A}\) 称为凸集 A 的度规。若 A 是吸收的且对称的，则 \(p_{A}\) 是一个半范数。

命题 23. --- 设 E 是拓扑向量空间。若 A 是包含 0 的开凸集，则 \(p_{A}\) 是有限的且连续的，且 \(\mathbf{A} = \mathbf{W}(p_{\mathbf{A}}, 1)\)。若 A 是包含 0 的闭凸集，则 \(p_{A}\) 是下半连续的，且 \(\mathbf{A} = \mathbf{V}(p_{\mathbf{A}}, 1)\)。

若 \(\mathbf{A}\) 是开的且包含 0，则它是吸收的。对 \(x \in \mathbf{A}\)，存在 \(\rho < 1\) 使得 \(x / \rho \in \mathbf{A}\)，从而 \(p_{\mathbf{A}}(x) < 1\)；这与 (9) 结合便给出 \(\mathbf{A} = \mathbf{W}(p_{\mathbf{A}}, 1)\)。由于凸函数 \(p_{\mathbf{A}}\) 在开集 \(\mathbf{A}\) 中有上界，它在 E 中连续（II, 第 18 页, prop. 21）。

设 A 是闭的且包含 0。对每个 \(x \in E\)（满足 \(p_{\mathbf{A}}(x) \leqslant 1\)），我们有 \(x \in \rho A\)（对所有 \(\rho > 1\)），因此由于 A 是闭的，\(x \in A\)；结合 (9)，这表明 \(\mathbf{A} = \mathrm{V}(p_{\mathbf{A}}, 1)\)。于是对所有 \(\mu > 0\)，\(\mu A\) 是诸 \(x\)（满足 \(p_{\mathbf{A}}(x) \leqslant \mu\)）组成的集；由于在 E 中 \(p_{\mathbf{A}}(x) \geqslant 0\)，这表明 \(p_{\mathbf{A}}\) 在 E 中下半连续（GT, IV, § 6.2）。

正次线性函数 \(p\)（\(\mathbf{E}\) 上的）是每个凸集 \(\mathbf{A}\)（其中 \(\mathrm{W}(p,1)\subset \mathrm{A}\subset \mathrm{V}(p,1)\)）的度规。

